% !PRINTED 273
\weberpage{273}
Fassen wir die Gruppe $L_p$ als Gruppe der linearen Substitutionen
\[
\eta = \frac{a\xi+b}{c\xi+d}
\]
auf, und lassen darin $\xi, \eta$ nach \S.\,65 die Zahlen $\infty, 0, 1, \ldots, p-1$ durchlaufen, so giebt uns $G$ die Gruppe der ganzen linearen Substitutionen
\[
\eta = a^2\xi + ab.
\]
Das sind die Substitutionen, die die Ziffer $\infty$ ungeändert lassen, und also eine Permutationsgruppe von nur $p$ Ziffern liefern. Es sind dies dieselben speciellen metacyklischen Gruppen, die wir im siebzehnten Abschnitte des ersten Bandes betrachtet haben.

Die verschiedenen aus $G$ transformirten Gruppen lassen irgend einen der übrigen $p+1$ Indices ungeändert, und die Gesammtzahl dieser Gruppen ist $p+1$.

\begin{center}
\S.\,71.\\
\gesperrt{Divisoren der Gruppe $L_p$, deren Grad nicht durch $p$}\\
\gesperrt{theilbar ist.}
\end{center}

Um die übrigen Theiler von $L_p$ zu finden, deren Grad nicht durch $p$ theilbar ist, können wir Schritt für Schritt denselben Weg gehen, der uns im siebenten und achten Abschnitte zu der Bestimmung der Polyëdergruppen geführt hat. Wir können uns hier damit begnügen, die Hauptmomente der Ableitung hervorzuheben und wegen der Beweise auf die genannte Stelle zu verweisen, wo ganz dieselben Schlüsse zu machen waren.

Es ist hierzu erforderlich, die Gruppe $L_p$, wie im \S.\,65 auseinander gesetzt, als Gruppe linearer gebrochener Substitutionen
\weq{1}{\Theta(\xi) = \frac{a\xi+b}{c\xi+d}}
aufzufassen, und darin der Veränderlichen $\xi$ alle $p^2+1$ Zahlenwerthe des Congruenzkörpers $\mathfrak{E}_{2,p}$, einschliesslich $\infty$, beizulegen. Dadurch gewinnen wir den Vortheil, dass die Gleichung
\weq{2}{\xi = \frac{a\xi+b}{c\xi+d}}
%CONT>>
